\section{分层码、联合交换与局部符号构造}
\label{sec:codes}

这一组七个结果来自同一个几何构造，改进却发生在不同层次。32、33、34维增加了
可以同时使用的八元支持；39维将更大的支持族与显式符号码结合；35、36、37维
进一步改变支持内部的符号选择。研究中的关键转折是：最初选择一整束点是否保留，
后来则允许用另一种有限点集替换这束点。

\subsection{分层构造的相容条件}

Edel--Rains--Sloane 构造\cite{ers}把三种坐标模式放到同一单位球面上。
设 $n\ge m\ge32$，$T\subset\{0,1\}^m$ 是二进制码，
$\mathcal B$ 是 $\{0,\ldots,n-1\}$ 的八元子集族。支持集记录向量哪些坐标非零，
其上的符号选择决定具体向量。稠密层、中间层和根层分别为
\begin{align*}
X_T&=\left\{\frac{1}{\sqrt m}
 \bigl((-1)^{c_0},\ldots,(-1)^{c_{m-1}},0,\ldots,0\bigr):c\in T\right\},\\
X_{\mathcal B}&=\left\{\frac{1}{\sqrt8}\sum_{i\in B}\varepsilon_i e_i:
 B\in\mathcal B,\ \varepsilon_i\in\{-1,1\},\ \prod_{i\in B}\varepsilon_i=1\right\},\\
X_R&=\left\{\frac{\pm e_i\pm e_j}{\sqrt2}:0\le i<j<n\right\}.
\end{align*}
每个八元支持有 $2^7=128$ 种偶负号模式。三层的非零坐标数分别是 $m$、8、2，
所以互不重合，合并后的点数为
\begin{equation}
N=|T|+128|\mathcal B|+2n(n-1).
\label{eq:ers-count}
\end{equation}

\begin{proposition}[三层构造~\cite{ers}]
\label{prop:ers}
若 $T$ 的最小汉明距离至少为 $\lceil m/4\rceil$，且不同支持满足
$|B\cap B'|\le4$，则上述三层的并集是 $n$ 维接吻构型。
\end{proposition}
\begin{proof}
稠密层中两个码字给出的内积是 $1-2d_H(c,c')/m\le1/2$。
同一支持上的两个不同偶符号模式至少有两位不同，内积不超过 $1-4/8=1/2$；
不同支持之间的内积不超过 $|B\cap B'|/8\le1/2$。
根层的不同点或至多共享一个同号坐标，或支持相同而至少一个符号相反，
内积也不超过 $1/2$。

剩下三种跨层点对。稠密层与中间层的内积至多为
$8/(\sqrt m\sqrt8)=\sqrt{8/m}\le1/2$；稠密层与根层至多为
$2/(\sqrt m\sqrt2)=\sqrt{2/m}\le1/2$；中间层与根层至多为
$2/(\sqrt8\sqrt2)=1/2$。这就覆盖了全部不同点对。
\end{proof}

于是，大型球面构型的几何条件归结为两个有限问题：寻找足够大的二进制码，以及
寻找两两交数不超过四的八元支持族。支持的指示向量具有常重八，并满足
$d_H(1_B,1_{B'})=16-2|B\cap B'|$；交数条件正是常重码的最小距离至少为八。
这一联系使常重码的改善可以通过式\eqref{eq:ers-count}直接转化为接吻数下界
\cite{bsss,brouwer,echols}。

\subsection{支持族的联合交换}

加入一个新支持，通常必须先删除与它交数过大的旧支持。若只考察单个候选，
删除代价可能抵消全部收益；几个相容候选却可能受到同一批旧支持阻挡。
研究因此转向同时选择新增支持，而不是逐个尝试插入。

对候选 $b$ 定义 $F(b)=\{B\in\mathcal B:|B\cap b|>4\}$。
若 $Y$ 与旧族不交，且内部交数均不超过四，则
\begin{equation}
\begin{split}
\mathcal B'&=(\mathcal B\setminus F(Y))\cup Y,\qquad
 F(Y)=\bigcup_{b\in Y}F(b),\\
|\mathcal B'|-|\mathcal B|&=|Y|-|F(Y)|.
\end{split}
\label{eq:exchange}
\end{equation}
收益取决于冲突集合的并，而不是各候选冲突数之和。所保存的33维最后一次交换
删去九个支持、加入十个；34维删去六个、加入七个；37维的一次交换删去37个、加入40个。
数据同时保存父支持族和实际删除、插入集合，可以重新计算这个并集。

这一目标还有一个直接的离散表达。给定有限候选集 $\mathcal C$，以
$z_b\in\{0,1\}$ 表示选择候选 $b$，以 $u_B\in\{0,1\}$ 表示删除旧支持 $B$。
对不相容的候选对要求 $z_b+z_{b'}\le1$，对 $B\in F(b)$ 要求 $z_b\le u_B$，
并最大化
\[
\sum_{b\in\mathcal C}z_b-\sum_{B\in\mathcal B}u_B.
\]
任何最优选择都只删除实际被所选候选阻挡的旧支持，所以这个目标恰好等于
式\eqref{eq:exchange}的净收益。它揭示了搜索的关键：应让候选的冲突集中在
共同的旧支持上，而不只是分别寻找冲突最少的候选。候选生成和组合选择是两个
相互配合的步骤；改变候选族，才会改变这一有限优化问题能够实现的构造。

阻挡集合也给出了候选生成的组织方式。给定整数 $f$，考虑旧族之外满足
$|F(b)|\le f$ 的八元子集。若候选池包含全部这样的子集，就覆盖了至多加入
$f+1$ 个支持的所有正增交换：对 $b\in Y$，有
$|F(b)|\le|F(Y)|<|Y|\le f+1$。这个阈值刻画的是具体的交换邻域，
而不是最终码本大小的限制。

候选之间的冲突还有一个紧凑表达。两个八元支持交数至少为五，当且仅当
它们包含同一个五元子集。每个候选占用 $\binom85=56$ 个五元子集，因此要求
\[
\sum_{b\in\mathcal C:\ Q\subset b}z_b\le1
\qquad\left(Q\in\binom{\{0,\ldots,n-1\}}5\right).
\]
这些约束可以逐批加入：求解当前选择问题，检查选中支持包含的五元子集，
再加入重复子集对应的容量约束。没有重复时，所选支持内部相容；若目标为正，
就得到一个有效交换。所保存的34维计算使用10715个至多有三个阻挡的候选，
通过这一过程得到上述七换六的交换。每次接受交换后重新计算阻挡集合，
又可继续探索新的替换。

最终直接用于三层构造的参数如下。37维的支持数先从2843增加到2845，
还将在下一小节作符号替换，故不列入这张直接计数表。
\begin{center}
\begin{minipage}{\linewidth}\centering
\captionof{table}{直接采用三层构造的参数。}
\label{tab:code-parameters}
\begin{tabular}{rrrrr}
\toprule
维数 $n$ & 稠密层长度 $m$ & $|T|$ & $|\mathcal B|$ & 总点数\\
\midrule
32&32&131072&1676&347584\\
33&32&131072&1803&363968\\
34&32&131072&1960&384196\\
39&39&327680&3383&763668\\
\bottomrule
\end{tabular}
\end{minipage}
\end{center}
每个新增支持贡献128个球面点。例如32维的点数是
$131072+128\cdot1676+1984=347584$。
33、34维相对于比较构造分别多15、24个支持，对应1920、3072个点。

长度32的稠密层来自 Cheng--Sloane 的 $[32,17,8]$ 码\cite{chengsloane}，
生成矩阵展开为 $2^{17}=131072$ 个码字。39维采用
$T=\bigcup_{a=1}^{10}(H+r_a)$，其中 $H$ 的维数为15、最小重量为12。
枚举所给核与代表，全部45对不同陪集满足
\[
\min_{h\in H}\operatorname{wt}(h+r_a+r_b)=10.
\]
所以它们互不相交，合计 $10\cdot2^{15}=327680$ 个码字，最小距离为10，
满足39维所需条件。核由扩展二次剩余码缩短得到；配套程序比较生成矩阵的行空间，
并直接检查实际码字的距离。

39维所列公开比较构造的稠密层规模相同，支持数从3324增加到3383，
故增加 $128(3383-3324)=7552$ 个点。将支持层和稠密层分别记录，
可以准确区分构造所用的有限对象及各自的贡献。

\subsection{Hadamard转移与局部替换}

继续增加支持数，并不是唯一可用的变化。一个八元支持原本携带全部128种
偶符号模式；若删去少量完整符号束，改用另一个带符号点集，局部区域可能容纳
更多点。关键是同时控制替换集内部，以及它与保留部分之间的内积。

\begin{proposition}[配对 Hadamard 转移]
\label{prop:hadamard}
设四元支持族 $\mathcal A$ 两两交数至多二，且相对于给定坐标配对均为横截支持。
为各支持附上全部符号，再在每对坐标上作用 Hadamard 矩阵，得到
$16|\mathcal A|$ 个平方范数为八、不同向量内积至多四的权八向量。
\end{proposition}
\begin{proof}
取两两交数不超过二的四元支持，为每个支持附上全部16种符号。将坐标两两配对，
只保留在每一对中至多使用一个坐标的四元支持，称其为横截支持。在每对坐标上作用
\[
H=\begin{pmatrix}1&1\\1&-1\end{pmatrix}.
\]
$(\varepsilon,0)$ 或 $(0,\varepsilon)$ 都变为两个非零坐标，
故权四向量变成权八向量。不同输入向量的内积至多为二：异支持由交数控制，
同支持的不同符号模式至少有一位改变。因为 $H^{\mathsf T}H=2I$，
变换后的平方范数为八、不同向量的内积至多为四。变换可逆，点数保持不变。
\end{proof}

将这些向量除以 $\sqrt8$，得到内部相容的球面点集 $Q$。
再检查每个新支持与每个保留支持的交数不超过四，即保证它们与保留中间层相容。
与稠密层、根层的交叉界仍分别为 $\sqrt{8/32}=1/2$ 和 $1/2$，
不依赖新点的符号奇偶性。这正是局部替换能够脱离完整偶符号束的原因。

自由区域内的局部替换原理已见于 Takhanov--Yun 对带符号权四码的研究
\cite[第V节，定义5与引理4]{signed-johnson}。这里将这一原理用于 ERS 中间层的
权八区域，通过区域选择、横截支持与配对变换实现下述具体替换。
35、36维中，边界相容性可以在选择 $Q$ 之前统一保证。取坐标区域 $W$，
删去完全包含在 $W$ 中的旧支持，要求每个保留支持 $B$ 都满足 $|B\cap W|\le4$。
于是，任何放在 $W$ 中的权八新向量都与保留中间层相容，搜索只剩局部码的内部
点对条件。35维可用 $W=\{12,\ldots,22\}$，36维可用
$W=\{12,\ldots,23\}$，此处坐标从零编号。它们分别包含一个和三个旧支持。
这样的区域把一个与整个大构型有关的问题化成独立的小型构造问题；删除费用分别
固定为128和384点。

直接搜索十一和十二坐标中的带符号权八向量，候选规模分别为
$2^8\binom{11}{8}=42240$ 与 $2^8\binom{12}{8}=126720$。
权四转移将问题改写为组合源与坐标配对的选择，同时保留内积条件的代数证明。
35维从十点 Steiner 四元系 $S(3,4,10)$ 出发：
每个三元组恰属于一个四元块，故共有 $\binom{10}{3}/4=30$ 个块，
不同块的交数至多为二。对十个坐标的945种完美匹配逐一选择横截块，
得到16个可用块。36维采用 Kalbfleisch--Stanton 的51区组构造
\cite[定理7]{kalbfleisch-stanton}；Takhanov--Yun 也在
\cite[第V节，构造1]{signed-johnson}中给出了这一构造。
将十二个坐标分为两个六元集，在每一半中采用同一组
$K_6$ 的五个一因子。每种颜色的左右两条边有九种配对，产生45个跨区四元块；
再在两半各加入第一种颜色三条边的补集，得到另外六个块。

这51个块仍两两交数至多二。同色跨区块若只共享一条边，交数为二；异色块在
每一半至多共享一个端点。同半区的内部块交于二点，不同半区的内部块不交；
内部块与跨区块的交数至多为二。

\begin{proposition}[576点局部码]
\label{prop:local36}
从两个一因子分解中选取36个横截块，经配对 Hadamard 变换，得到十二坐标中的
576点权八码。
\end{proposition}
\begin{proof}
在两半中分别以同一颜色的三条边作为坐标配对。其余四种颜色的每条边，在每个
已选坐标对中至多取一个端点；左右两条同色边的并因此是横截块。
这样的块共有 $4\cdot3\cdot3=36$ 个，且由前述交数计算，两两相交至多二点。
命题~\ref{prop:hadamard}给出 $36\cdot16=576$ 个点。
六个内部块不参与这一横截子族，故存在性证明无需使用它们。
\end{proof}

在所给局部坐标顺序中，实际配对为
$(0,1),(2,4),(3,5),(6,7),(8,10),(9,11)$，恰好是两个半区中的同一颜色。
采用这些配对的直接构造与保存的576个向量完全一致。
配套程序还枚举两种源码的945与10395种坐标匹配，重建256点和576点集合。
这样，点集的存在由显式构造证明，穷举计算则给出这一有限模型中的配对选择。

三份有限构造分别选择16、36、32个横截支持，产生256、576、512个新点。
35维的相容区域有11个坐标，构造将其中10个配成五对，余下一位恒为零；
36、37维各使用12个坐标。配对与横截支持均保存为显式生成数据，
可以重新生成256、576、512个带符号点并与最终点集逐项比较。
\begin{center}
\begin{minipage}{\linewidth}\centering
\captionof{table}{局部符号替换；原支持数指删除符号束之前的支持数。}
\label{tab:signed-replacements}
\begin{tabular}{rrrrr}
\toprule
维数 & 原支持数 & 删除的128点束 & 加入点数 & 替换净增\\
\midrule
35&2157&1&256&128\\
36&2742&3&576&192\\
37&2845&3&512&128\\
\bottomrule
\end{tabular}
\end{minipage}
\end{center}
合并支持改善与符号替换，得到
\begin{align*}
K(35)&\ge131072+128(2157-1)+256+2(35)(34)=409676,\\
K(36)&\ge131072+128(2742-3)+576+2(36)(35)=484760,\\
K(37)&\ge131072+128(2845-3)+512+2(37)(36)=498024.
\end{align*}
37维相对于2832个支持的比较构造，先由新增13个支持得到1664点，
再由局部替换得到128点，合计增加1792点。35、36维保持原支持数，
全部增益都来自符号选择的变化。

35维的局部选择还有一个恰好达到等号的容量结论。令 $M(11)$ 表示十一坐标中
权八的 $\{0,\pm1\}$ 向量集在不同向量内积至多四时的最大大小。给定完整符号串
$\sigma\in\{\pm1\}^{11}$，与它在非零位置符号一致的两个权八向量，
支持至少交于五个位置，故内积至少为五。因此每个完整符号串至多容纳一个向量。
而每个向量在三个零坐标上有八种符号补全，计数得到
\[
8|Q|\le2^{11},\qquad M(11)\le256.
\]
上述256点构造达到等号，故 $M(11)=256$。这个局部等号解释了为何35维的
十一坐标区域能够被完全利用。十二坐标中两个八元支持可以只交于四个位置，
同一符号串便可能容纳多个向量；这正是36维能够使用更大局部码的几何原因。

有限检查使用数据明确给出的坐标顺序、配对和四元支持。37维支持串从左到右
以一为起点编号时，坐标 $i$ 对应整数掩码的第 $37-i$ 位。
程序据32个横截支持和六对坐标重新生成全部512点，检查与保存的替换集一致，
再检查与所有保留支持的交数。构造的来源、点数和几何条件由同一组输入连接起来。

这七项结果将支持族与其上的符号模式作为两个可以共同设计的对象。
下一章转向格构造：其中可改变的是大型提升构型尚未使用的赤道方向。
